\section{The Establishment of a Traditional Society: Priests}

\begin{quotex}
In \textit{La Tradizione Romana}\footnote{\textit{The Roman Tradition}}, \textbf{Guido De Giorgio} dedicates a section to the question of the reestablishment of a traditional society. It includes four chapters: \emph{The Priests}, \emph{The Warriors}, \emph{The Workers}, and \emph{The King}. This is the first part of the chapter on the priests. De Giorgio provides some very precise indications which are very demanding and ought to be taken with some seriousness.
\end{quotex}

\subsection*{Introduction}

The general lines that are given here about a society established along the norms of a tradition truly such, can serve as the model and example for leading the current West to a normality that it has lacked for quite some time, which could, however, be gradually and prudently restored without great difficulties by the return to those principles and to that spirituality of thought and life that are indispensable to the very existence of men. Therefore it is a question of a general type according to which Europe could be oriented and hence the world with the necessary adaptations to change conditions of life while avoiding a violent and radical dissolution that would jeopardize, through its rapidity, the reintegration of a normal state.

If a return is possible, if a rectification is still possible to save the world from spiritual and material ruin, it entails a change of direction that must absolutely proceed from the inner to the outer and not otherwise. First of all, spiritual conviction of a similar change is necessary, and this conviction must be founded on a \emph{purely intellectual realization}, i.e., without any sentimental infiltration whatsoever, of principles that rule a traditional society, and, since it would be useless and childishly utopian to believe that the entire people and the masses can achieve it through an instant process, it would be sufficient that a few at the beginning would constitute a formative nucleus with an increasingly more efficient operative irradiation leading up to the reintegration of a nearly perfect form. These few would be animated by a single force, that of truth, and guided by a single goal, that of making it triumph, and not to give in to any concession toward oneself of others, but to become true and proper centers of salvation for men misled and corrupted by centuries-old errors. If such men do exist, the return to normality would again be possible and the truth could still triumph over ignorance.

A change of this type entails the pure and simple rejection of all the prejudices that have debased Europe for centuries, of all the errors that the Anti-tradition has been amassing for hundreds of years to corrupt thought and life, impoverishing the mind and paralyzing those spiritual forces promoting the true and only good, the fruition of the Truth in a superhuman order where it alone resides and where it is necessary to have the courage and the strength to maintain it. These men could save Europe and the world, leading the people back into the great traditional channel, toward the light of those principles that are the very basis of existence. But a great energy is necessary to conquer pessimism and skepticism which are opposed to any reestablishment on the one hand, and superficial optimism, on the other, that considers apically reached that which is a simple degree of transition. Above all, great detachment is necessary, strict intellectuality, the absolute absence of that pseudo-mysticism so much in style at the present time where spirituality and sentimentality are equivalent, where enthusiasm and faith are placed on the same level, where the darkest and murkiest impulsiveness is believed to be an expression of strength, where the external is not only excessive, but tends to destroy the internal, where agenda kills true development, where finally all that is inferior and illegitimate is affirmed with an immodesty which the world has never known before now even in the most acute periods of decadence.

But above all, great courage would be necessary in these men and a great faith to conquer so much darkness of ignorance, to break up so many false bridges, to facilitate the return to the comprehension of the truth, to give back to thought its dignity and to life, its justification. Many currently in Europe are aware of the mud the world gets mired in every day, but their conviction is half-hearted and imperfect: it is more about a pessimistic and skeptical attitude than a true and proper conviction. None of those would be able to indicate the remedy or to show the way to return to true life, which, before all else, reconfirms the values of the spirit that are divine, but do not neglect the necessities of existence, orienting them along the purely traditional axis.

Every traditional society entails the organic allocation of labor and the assignment of tasks to various categories of men adapted to their nature and their possibilities. Even currently, in the agonal state of Europe and the world, this allocation exists, but it is false, arbitrary, unnatural, temporary, not founded on traditional norms, i.e., according to a strictly and purely hierarchical order that derives from the very meaning of the term \emph{hierarkhia} = high priest in Greek, whose sacred character should not elude anyone. Since every tradition is sacred, the allocation of labor must be imposed only from deference to the truth that is the divine order; therefore, the priests would stand at the apex of a traditional society, who are the bearers of the divine order, the divine knowledge. We could say Ascetics, but we prefer to preceding term because it better highlights the precise character of the nature and the office assigned to the bearers of supernatural truth. This caste includes those who do not participate in active life, who lack every responsibility of the profane, temporal order and are the leaders of traditional society, but in reality, through the very character of their mission, occupy a marginal place in respect to the practical act of existences, but one that is central, determinative and indicative for the maintenance and the development of traditional society. They are the poor of God, the voluntary renouncers of the world, the truth givers, the Spirit lovers, the last because the first, whose life is devoted to the realization of the great divine norms that constitute the traditional body. No confusion is possible regarding their nature and their missions which are of a purely spiritual order, without any worldliness, extending into the pure regions of the Spirit of God who governs the world invisibly.

\subsection*{Process of Degeneration}
\begin{quotex}
This is the second installment from \textbf{Guido De Giorgio}`s chapter on the priestly caste from \textit{La Tradizione Romana} \textbf{The Roman Tradition}. He sees the process of the degeneration of the castes as beginning, not from the revolt of the Kshatriya, but rather from the failure of the Priests to fulfill their divine mission. Nevertheless, Tradition may be forgotten, but never lost, so there are always some solitary men who continue to maintain that sacred knowledge. De Giorgio moves on to the topic of faith which, he says, is necessary because a man must believe before he knows.
\end{quotex}

Therefore, they constitute the caste of the invisible ones because their continuous, untiring, hidden action is of the sacred spiritual order and is spontaneously fulfilled through the very force of the soul, the purity of their life that must be a rite, an offering, a sacrifice. We insist particularly on the character of the priestly caste to determine the modality of the hierarchic power and the efficacy of the revelatory mission entrusted to them. They are the leaders of traditional society, but invisible leaders because the sacred knowledge of which they are the depositaries is not a purely dogmatic system nor a dead body, but a living, perennial fire that they must nourish, living continuously in communion with the spirit of God, acting so that His truth is towering at its peak, in isolation from His profound virtue, in the active and effective development that only the contemplative life allows to be realized.

They are the great hermits and they descend among men, the bearers of the graces of the truth, but, even if living in the world, they are really on the outside and dominate while not dominating, act while not acting, illuminating with their light, they save with their presence, strengthen with their example. Their hierarchy and their organization, the traditional type to which they are tied, the value of the character of their mission are absolutely determined by their having freely and consciously chosen not a career, not a profession, not a job, but the way of God, that by leaving from God leads back to God, and of having chosen it not only for itself, but rather with the precise task of showing it to all. They sacrifice themselves and they sacrifice: it is reflected carefully in the value and in the immense importance of these two expressions when led back to the precision of their etymological meaning sacrifice=make sacred. The Priests cannot not be sacred themselves, they cannot not make sacred everything that they touch and go near because they were born for that, they were destined to that and their choice, in embracing the holy ministry, is an absolute match with the inherent possibilities of their nature.

They cannot deviate if they are truly priests, they cannot, if they deviate, not fall into the lowest abjection because they have betrayed God, they failed in their vows, and by abjuring, they contaminated themselves and men. If the value and the height of the task entrusted to the Priests is fixed and well understand, one will reach some conclusions to which the modern world, having become anti-traditional, is no longer capable of elevating itself. Let us outline one of them.

Since the Priests are the holders of sacred knowledge and constitute the insuppressible base of a truly traditional foundation, they are to secure its normality, they are to maintain the unity with their invisible spiritual hidden work and therefore are responsible for the general defection of the traditional spirit, because no one can fall if the priests themselves do not fall, no one can fail if the priests themselves do not fail, no one can be banned against the truth of God if the priests do not betray it first, no one can corrupt the world if they do not corrupt it first, abandoning their sacred mission for the concerns of the temporal order, passing from the contemplative life for which they are destined, to the active life that is not absolutely the place of the development of their activity, with falling short of their caste, their obligations, and especially the divine principles whose radiating virtue they must keep intact. How many are capable of understanding how \emph{the current state of abjection is due to the defection of the priestly caste who are responsible for it} because they alone, the Priests, maintaining contact with the divine not only by means of their action over men, but above all with the constant realization and efficacy of their interior ascesis? Just as God's strength is mysterious and invisible, so is the Priest's occult and hidden: contemplating they act, fulfilled in God, they work in the world, sacrificing themselves they sacrifice, praying they save, although their mission is genuine and not the unholy profanation of God's laws.

If one can with absolute certainty impute to them the state of current Europe, whatever defection of this caste, whatever decadence of humanity is not attributable to the traditional form to which they are reconnected and of which they could be the authentic representatives. Tradition is invulnerable, inviolable, unassailable, it is God's truth and is kept intact because, even if betrayed by His ministers, He always finds those who preserve its sacred character, those who, among men not belonging to the priestly caste, become its legitimate and authorized carrier. And almost always those, priests among men, carry out their mission more dangerously than if they belonged to the officially recognized caste, because they have to fight against a profane force that tends to suppress them, that of those who have betrayed the faith, repudiated tradition, abandoning the divine for the human and the sacred for the profane: those are the false priests who are no longer such, i.e., holders of sacred knowledge. Reflect on the importance of what we say, and you will be able to understand how the solitary Ascetics, those who we will be able to call ``hoi exo'' the outside , in all the periods of decadence of the priestly caste, have kept alive the perennial fire of tradition against the deceit, the hatred, calumny of those who failed in their mission.

The self-establishment of ascetic groups outside the priestly caste, the presence of Masters, i.e., of solitary Ascetics, in all the period of decadence is explained precisely by the abandonment of tradition by parts of those to whom its deposit was entrusted, and the necessity---of the divine order, we insist---that others seek to keep contact between man and God, purging the sacred ways of the profane dross that the false priests amassed, i.e., the most ungodly deniers of the supernatural world. Dante \emph{docet}... It is necessary, to avoid misunderstanding, to insist on the nature and character of this definition, of this contamination that takes place in the ambit of the priestly caste in the periods of decadence. The divine truths that constitute the sacred body of tradition has a purely metaphysical, or transcendent, character: they are superhuman, eternal, and to approach them, it is therefore absolutely necessary to pass beyond the human condition and bring oneself with the intellect into that sphere of pure actuality where the divine reality is developed beyond the domain of Forms and Rhythms in the silence of its ineffability. Faith prepares this transition from the human to the divine, in fact, it is its essential condition, that which cannot elude anyone through the elementary analogy with so many human and contingent situations. It is necessary to believe that one knows, because one knows only by knowing that one can one know before knowing, i.e., before having acquired the wisdom and having already completed the passage from the human to the divine in order to do so that only the divine is.

Faith in whom? In God, in the Master, so say all traditions that insist on this absolutely necessary condition for the effective realization of the divine. One believes in the truth before reaching it, i.e., before being there and being it, and the intensity of faith is in direct relation with the efficacy of the achievement.

Faith is therefore the rail, the bridge, the isthmus between the human and the divine, between what man is not and what he really is when he is no longer, when he surpassed and passed over forever the human condition. But since this is the fruit of ignorance, faith is the necessary condition for the dispersal of ignorance and the reaching of wisdom.

It annuls every human limitation in man, abolishes individuality, opens all the passages of Infinite possibility, considers the chains as untied so that they are really so, it works a type of preparatory radiation of individual faculties, because one believes in things other than oneself, in the sacred text, in the hidden value of the rite, in the minister, in the master, in other words, in whatever passes beyond quotidian reality, the illusion of the world ordinarily experienced in the ambit of all sensible and rational limitations, it denies resolutely the tangible whole and affirms an invisible reality.

To have faith means to believe in what one does not yet know, one does not know, it is the most noble and desperate attempt to bring oneself face to face with the threshold of mystery and to affirm that beyond it there is an unspeakable reality, that which is revealed. Even for those who cannot pass over this threshold, it is enormously positive that they succeed in bringing themselves to the extreme limit connected to their strengths with an act of faith that leaps over error, the visible presence, the world and things, in one stroke, to genuflect in the face of the Invisible Presence, God and His Names still concealed precisely because they are absent, asserted, and believed and not known. Faith is therefore superior to any human knowledge whatsoever, to any conquering activity of the human alone, because it passes beyond it, considers it as ancillary, negligible, negotiable, even nothing in the face of the divine that it recognizes as the invisible venerable root through its invisibility, real through its apparent unreality, divine exactly because not human, not tangible, not discernible with the senses, and analyzable by reason, placed in a sphere of this the good ones, those who believe, will be able to enjoy, if it will be pleasing to God, only after the dissolution of the human compound, i.e., \emph{post mortem}.

Faith works this miracle that those who cannot reach by an operational effort and aware of the threshold of the divine, attain it with the rapid and direct contact that can be fecund with greater results: compared to those who know, that are beyond the threshold, they find themselves very far, but compared to those who do not believe, to the small men of the small world, the affirmers of the minimum in the minimum and lovers of the shadows, they are in a clearly privileged position because they are distant from just as far as the spirit surpasses the flesh and intelligence surpasses imbecility.

\subsection*{Conclusion}

\begin{quotex}
This is the third and final installment from \textbf{Guido De Giorgio}`s chapter on the priestly caste from La Tradizione Romana \textbf{The Roman Tradition}. He continues the discussion on faith. He points out that when the priesthood deteriorates, it gives rise to priesthood of solitary Ascetics. Although they maintain traditional teachings, these solitary men do not have the authority to maintain traditional society. Furthermore, they find themselves opposed by the decadent priests. The priests who go astray need to be punished in order to prevent scandal. In the contemplative life, one's own thoughts and desires are emptied and the Invisible Presence begins to manifest in one's center in an endless number of stages. By losing his life in this way, his life becomes more real and more fulfilled.
\end{quotex}


Those who are called ``geniuses'' by men of today, so tied to the human to see them through the greatest irony, even the superhuman, enclosed as they are in human and terrestrial limitations, are greatly inferior to the most humble of believers because they hypertrophy a nothing, man and the world, and they make everything from this nothing, while the believers deny the nothing, man and the world, and reaffirm it only putting it back on God, i.e., in the supreme cause. Those who do not believe see the effect separated from the cause, which is absurd, while those who believe see the effect in the cause, which is conformed to truth: those who know even abolish the effect and this is the truth. Deepening this last point, one will be able to come to understand, first of all, the triple attitude of man in the face of the truth in accordance to what is subhuman, human or superhuman, and consequently the way in which the principle of causality must be considered in the transition from the profane to the sacred, from the human to the divine, according to what the schema of death, life, and liberation makes of it. This triple schema can be more crudely formulated in connection to these possibilities: \emph{without God, with God, God}.

Faith is therefore the traditional base for excellence since it is the necessary anticipation of the means through which the path begins from the human to the divine and it resolutely builds a bridge that is fixed on the other shore without seeing it but only knowing it as revealed: tradition works this fixation in the divine that constitutes for the believers the invisible \emph{ubi consistam} place to stand , the fulcrum of its elevation to the threshold of mystery.

We note that God, precisely because He is believed but not known, is affirmed in His deepest reality, that of the Unmanifested Principle, and that simple faith pronounces itself more positively than what would appear to a superficial examination, since by affirming the unknowability of the destination, it admits implicitly that only a realizing knowledge can attain it with an effective becoming. This knowledge follows faith, it passes over the threshold of the divine to carry oneself into the same mystery.

It is the most interior part of the Temple where the sacrifice is performed, that altar where the Priest resides.

So if the whole temple constitutes the domain of the sacred and we have the profane only outside of it, in the temple itself among the altar and the rest, there is a clean separation, and while the altar constitutes the active realization of the divine and therefore it is the true domain of the sacred, the rest of the church is reserved to those who are present at the sacrifice, they participate in it, \emph{but do not perform it themselves}, therefore they are always only the profane because there is a barrier between them and the divine truth. This is the relationship between faith and knowledge, between the faithful and the Priests, between those who only believe and those who know and should know, between the throng of the \emph{called} and the and the rank of the \emph{elect}, between passive love that genuflects and pulls back outside of the sacred threshold and the active love that performs the sacrifice directly with a gesture that is a benediction, a voice that is of God, an altar that is the very throne of the Eternal.

From all that it is easy to reach this conclusion that the Priests, as the bearers of sacred knowledge, must really possess it by means of a realizing knowledge in order to be able to perform on the altar the highest of the sacrifices, the divine sacrifice, through which they strengthen the faith, reconfirm it in its strength, and reach the hope that it can, by crossing the threshold, reach the sphere of
\emph{charitas}, of divine love.


If the Priests do not attain this realizing knowledge, they are not such, they cannot be the sustainers of Tradition, the princes of a traditional society, those who maintain the contact between the human and the divine in a permanent current, who secure life and the justification of life itself, placing it back in God. This is the true sin which a Priest can incur, of not having attained the realizing knowledge, of performing a rite with knowing its value, efficacy, and meaning, of not knowing, in other words, the sacred knowledge and of reducing it to a simple babbling of the lips or to an idle formulation that, in the scope of realizing knowledge which should be that of the Priest, is an absurdity because the one falling short puts himself beneath the believers, the profane among the profane and the profaner among the profaners.

Tradition remains untouched however, even in such a case, the intangible rite does not diminish the efficacy that it had on believers, on the faithful, some of whom can even substitute, if not in fact, in a reality that even leaps over the fact, the Priest, performing for them the rite and becoming for them the transient or permanent deposits of God's truth. If this defection does not harm Tradition in itself, it harms enormously the basis of a traditional society, it invalidates the principle of authority that is absolutely necessary to the maintenance of the two orders, the human and the divine, the profane and the sacred, that of faith and of knowledge, of those who believe and those who know, but what is more serious, it gives rise to the \emph{priesthood of solitary Ascetics} who, by God's will, maintain the traditional secret with their work reserved to a minority, they safeguard it, they protect it against profanations, and they end up being opposed by the priestly class that fears in then its same knowledge, that which they themselves abandoned and betrayed. These deplorable clashes always occurred at the expense of sacred knowledge, giving rise to upheavals and conflicts of every type contrary to the maintenance of traditional order.

In making the Priests the depositaries of sacred science in in considering so these various questions, we thinking of a traditional form that more is adapted to nature and the spirit of the West where the sole ascetics should be the Priests, where they along should truly safeguard the traditional body realizing effectively in itself the doctrine with the deep knowledge of everything that constitutes the divine truth.

The Priests besides being deprived of real knowledge of sacred things can also come up short of the decorum that is imposed on their caste. The accusation of immorality is too noted and too usual because one has to insist there, instead of this were derived incomprehensible devaluations of tradition itself, which is truly absurd, giving rise to the sect, schism, heterodoxy. We will immediately and clearly say that the morality is for knowledge and not the inverse and the purely moral side of a question of any fact is the least act for make it comprehend and valued its meaning. The first and only sin is ignorance and one is responsible more for not knowing or poorly knowing than to act evil. As to judge conscience, we do not believe that his is not only a very arduous and insolvable task, but directly a sacrilege, because from God alone and in the name of God consciences are judged and where it never than the so called more judgments are given in God and in the name of God? One reflects attentively that and one will see that only the facts can be judged in the customary place as the condition capable of determining others and therefore, if harmful, condemnable.

Therefore what must be judged, condemned, and punished is the deed when it is visible, it is ``active'', it leads to imitation, it is the example, the seed of other deeds, it corrupts: what must be condemned is scandal in all forms, in all classes, and this is the healthy principle of a truly traditional ethic that maintains the privilege of the conscience in action as God's privilege, and represses the infraction not because the motives are condemned, but to prevent the repetition, the example.

And since there is a vast range in the intensity of the propagation of an offense, i.e., in the immortality of a deed, a truly civil society should adopt all those natural, external, brutal, repressive measures that punish, from the death penalty up to the torture and flogging, but
\emph{they do not judge}, in beating the man, what he has that is the most external, i.e., the flesh.

A truly civil society, readopting all the old measures of physical, outward coercion, returns to the truth of the traditional river-bed, obeys scrupulously God's truth, and at the same time enriches the number of true, real reasons for action, creating a game, an alternation, a reciprocity between the crime and the punishment, fruitful in sensations beneficial to life, fertile in expansions and the surpassing of the brutal fact. Capital punishment, the various types of torture publicly inflicted with their tragedy, are always considerably effective and instructive and they can even determine true sources of purification, positive complexes, whose importance does not elude those who are equipped with sense and truly constructive imagination and not esthetically and softly deviators. But there is another consideration of a deeper order: the crime has degrees integrally subordinated to human passionate nature, i.e. to the darkest complex that dominates ordinary men undisputedly: it is therefore necessary that hierarchically or analogously proportionate penalties correspond to these levels, penalties determinate by an impersonal justice that strikes in the flesh, abstracting from any experimental character: thus only a reward and a rectification reaching a balance is established; this is the indicator of a true and proper traditional society.

The Priests can and must be blamed only if they fail at their task through total or partial incomprehension of traditional truths: their existence, on the purely exterior side, i.e., from the moral point of view, is reprehensible when it offers matter for scandal and then only if they submit to a tribunal that can be constituted only by men belonging to their caste. In fact from what has been said thus far, it results that real knowledge, i.e., integrative and realizing, belongs to the contemplative, and not the active life. The relationship between contemplation and action is of the greatest importance for the maintenance of the traditional idea because their imbalance constitutes a rupture, a hierarchical inversion, a veritable deviation that, prolonged for centuries, was the origin of the current abjection of the West. Contemplation stands to action in the same relationship as the divine to the human, the sacred to the profane, the eternal to the ephemeral, since their range is distinct and cleanly circumscribed by the two different types of activity. In contemplation there is an activity of a special order that is accomplished in the eternal place, beyond time and space in the sphere of transcendent truths, in an apparent retreat and in an interiorization that in reality are a veritable translation from the human to the divine and a cancellation of the human so that only the divine remains in its absolute autonomy.

In this sense, contemplation and revelation are synonyms because the divine truth can be revealed only in man becoming a temple, i.e., the refuge itself of the truth, the empty temple of humanity founded on earth and elevated to heaven in a verticality symbolically reflecting the lifting up of all being to the totalization of the higher states for their inclusive and intensive integration.

If in the Primordial Tradition the world itself was the temple, with the successive degeneration of humanity, the sacred enclosure was established to separate the sacred from the profane and maintain the distinction between the two orders in the way that the higher directs and guides the lower. The temple is the symbol and more than a symbol, it is the place of peace, of absolute interiority where, every individuality is denied, annulled, or distanced from every human dross, the realization of the divine is accomplished, the theophanic cycle of all its effective fullness.

Whoever contemplates---and contemplation is only of the divine order, therefore it must be judged absolutely inappropriate to any order, especially to the esthetic order which is visibly inferior---not only separates himself from others, but from himself, that which is the essential, and it empties his heart making it the center of his being where the Invisible Presence is manifested in a progressive irradiation whose levels are without end and constitute the hierarchy of divine stages. This term therefore cannot be, nor must be, applied to another that is not the real, effective achievement of higher states not passively glimpsed as though from the outside, but actively realized in one's interiority of the great temple which is the purified heart, cleansed, made a receptacle of light, the holy chalice where the divine mystery is completed. Everything that is purely and clearly human---like art in the profane sense and philosophy especially in the modern meaning---is excluded from the contemplative field that is the divine and not human life, reality and not illusion, truth and not ignorance. Philosophy that is an impaired wisdom and art that is a purely exterior infatuation for their very degeneration, are excluded from contemplative life and represent an artificial interstructure that the current abjection established between contemplation and action, small spurious world where weakness and imbecility produced by the ignorance of external things is depleted.

If contemplation is therefore reserved to the Priests who are the bearers of divine wisdom, what will be the relationship between the contemplative and active life? It is identical to that which rules the divine order and the human order: the active life must be oriented in accordance with a vision that can be determined only by those who live contemplatively, In fact, if man and the world in themselves are nothing when separated from their cause and not put back in it, i.e., in God, they acquire instead quite another meaning when they are integrated in the real order because they represent the same place where one of the divine possibilities is fulfilled. Reflect attentively on this: the contemplative life is a larger circle that includes in it a smaller one, the active life, the integration of these two modes constitute the traditional unity. Affirming the superiority of the contemplative life, we postulate the necessity of the active life provided that the latter is included in it, i.e., it remains hierarchically subordinated in order to shape it and procure for it all its operative efficacy. One, contemplation, works in the eternal, the other, action, works in the ephemeral, which is the symbol of the eternal: one, contemplation, is the domain of the sacred, the other, action, is the domain of the profane that becomes sacred only if it receives light from the former. To be more explicit: the contemplative returning to life properly called, i.e., to external existence, sees it under another aspect, different from that of the common man, therefore it is logical and natural that he seeks to consider it not as separate from the divine truth, but as effective preparation for that, the vestibule, the pre-position that rules and determines a final development realized only in the place of pure contemplation.

The harmony between contemplation and action is necessary for the complete integration of human possibilities so that man could be truly such, rich in every development, judge of his own destiny, capable of elevating himself from earth to heaven in a progressive expansion of all his faculties. But that is possible only if contemplation dominates action: in the reverse case we have the hierarchical revolution, the annulment of the traditional axis, the impoverishment of man demasculinized, made ugly, victim of all those inferior forces on which he can have the upper hand only if he is guided by the spirit of God.

It is up to the Priests, besides, to act so that the spirit of God reigns in the world and sustains the forces of man and truly dignifies him because man is everything with God, nothing without God, and his action deprived of any traditional content is a groping in the lower darkness, a wicked shadow among wicked shadows, sterile bedroom activities \emph{fornicatoria} that impels him from illusion to illusion, from error to error in the great detrital riverbed, the hellish, subhuman, ghostly, and lemuric world.

But if instead the active life is regulated along the traditional axis, as a rite, then all the apparent imbalances are annulled while counterbalancing each other, all the varieties of human espression and deviations are arranged and are resolved in an integrative homogeneity. It is necessary to insist on what men have completely forgotten in the collapse of all traditional values: the active life is so much richer the more it is subordinated to the contemplative life, because in a truly traditional society there is an amplification of human possibilities, an infinitely more fertile development of their most various activities, in the modern world, they are completely wasted. For a law of analogy that regulates the parallelism of the two orders, we will say that the more intense the contemplative life is, the richer is the active life, as much as God is exalted, so much more is man valued, and that in a truly traditional society, everything profane is something sacred \emph{in fieri} pending and all aberrations are recomposed in the equilibrium between the divine and the human, love and hate, wisdom and ignorance, war and peace, virtue and vice, good and evil, in a full, integral, oceanic confluence, are harmonized, surpassed, and twinned, in the great traditional river-bed.

But because this takes place, it is necessary for the Priests to truly be the supreme club-bearers and that their life be purely contemplative, as the bearers of divine truths, the great fee men of the eternity that they keep watch on the bastions of time, intermediaries between the higher and lower waters, between the divine and the human, between the cosmic and the hyperuranium where the Platonic ideas reside, knowers of the light that is the only light emanating from God.


\begin{center}* * *\end{center}

\begin{footnotesize}
\begin{sffamily}
\texttt{Janus on 2013-02-10 at 21:53 said:}

This is essential, and something many ``political'' traditionalists should read. There is an unfortunate trend toward restoring the forms of Tradition before its essence. Seeking out a spiritual or religious discipline would probably be the most important thing for those of a contemplative (brahminic) nature to do at the current time. Providence may yet provide teachers in surprising places.

That said, the essence is gaining more ground in surprising places. While being aware of the somewhat ambivalent attitude of many on this board to Dugin, it would be worthwhile to look at the Against the Postmodern World conference held in Moscow some time ago. No national bolshevik rally this; Shaykh `Abd al Wâhid Pallavicini was among those who saw fit to attend. Behind the petty racialists, neopagans without understanding, new right and national revolutionaries seems to be a stirring of some who are ready to look beyond the political form and see the inner life it is meant to reflect, be they of a contemplative or active nature. I'm especially interested in seeing what follows here. For those interested, here is an article from Dugin's conference: it is by Claudio Mutti and entitled ``The Doctrine of Divine Unity in the Hellenic Tradition''. Cologero, I'd love to hear your thoughts on it.
\url{http://against-postmodern.org/node/127}


\hfill

\texttt{Logres on 2013-02-11 at 11:40 said:}

Iamblichus would say that ``courage'' or ``daring'' is the Dyad -- the warrior -- the feminine principle: hence, the purely ``political'' is the Dyad without the Monad, doomed to futility. Very lucid article!

\hfill

\texttt{David on 2013-02-11 at 22:44 said:}

«A Christianity that is largely without doctrine and sacrament is a Christianity of slogan and extravaganza. A ``Churchless'' Christianity is simply, a heresy. It is a strange reading of the New Testament with conclusions as novel as they are effective. It is also destructive of the long term health of the Christian faith. Many who grow tired of its slogans and extravaganza do not turn elsewhere -- they turn nowhere. The fastest growing religious group in America is the unchurched.

The truth and richness of the Christian faith is only found in the deep-woven fibers of the historic Church. The life of sacrament, rooted in a thoroughly Christianized network of families, parishes and monasteries, is the normative existence of the Christian faith. This is the faith that converted the Roman Empire and the barbarian ancestors of people like myself. From it grew a great civilization, one that has been challenged and dismantled at many points, but which has yet to disappear.»

-- Ft Stephen. \url{http://glory2godforallthings.com/}

I.e., the Church is Tradition. Priest are part of a tradition and we need it back.

This author is very interesting for those who are Christians.

\hfill

\texttt{Jason-Adam on 2013-02-12 at 14:46 said:}

Guido de Giorgio's book is amazing and I wish a full translation of it existed. Cologero, do you know of a link to the full text in Italian online ?

As for Dugin, I know the man personally and can say that he does not act like an aristocrat -- meaning that he does not discuss and debate logically and cordially, the way that Cologero, Avery, Logres and all the folks at Gornahoor do, insteadhe expects one to accept his words without question and obey them, if one raises an objection or merely asks for further clairification of what he means, he gets angry and then starts calling names like ``reactionary'' ``CIA agent'' and down the line...... I ultimately broke with the Dugin camp after I realised that, at least in the west, he is totally uninterested at all in monarchy, Catholicism, anything from the west's past, he only wants to talk to communists, neo-nazis and those ilk...... 

\hfill

\texttt{Senkosmos on 2013-02-12 at 21:00 said:}

I need to agree with all of you: Cologero, David, Logres and Janus. Priests are the backbone (so to speak) for the preserving of tradition, and an instrumental tool to the intellectual elite (if they are not the elite themselves) to put the West back in the right track.Today, most agents of tradition are more concerned with political action than individual/local action, and more important, contemplation, spiritual non-acting action. All attempts to a ``back to roots'' without a ``back to the Absolute source'' is futile and ridiculous. Great post!

\hfill

\texttt{caseydeann on 2013-02-12 at 23:13 said:}

Great comment, David. Thank you.

\texttt{Logres on 2013-02-18 at 11:07 said:}

This is an immensely penetrating article -- I can see why Guenon regarded him as a peer. It's interesting that, like Kukai (the Japanese master), he places those of ``faith'' far above the secularists and agnostics, as pilgrims on the ladder up to heaven who have begun the ascent, which will bear fruit in time, or out of it. Did Giorgio consider himself part of the priest caste?

\hfill

\texttt{Cologero on 2013-02-18 at 12:26 said:}

I suspect he sees himself this way, as he described:

\begin{quote}

Reflect on the importance of what we say, and you will be able to understand how the \textbf{solitary Ascetics}, those who we will be able to call ``hoi exo'' the \emph{outside} , in all the periods of decadence of the priestly caste, have kept alive the perennial fire of tradition against the deceit, the hatred, calumny of those who failed in their mission.
\end{quote}


Certainly, he understands the nature of the priestly caste. When we get to the warrior caste, we will see that he describes them in terms very much like Evola. However, Evola was tone deaf to the contemplative caste and never understood it or its true role. That is why he saw the Ascetic and the Warrior on the same level. De Giorgio, on the other hand, will make their real relationship very clear.

\hfill

\texttt{Tosti on 2013-02-18 at 14:59 said:}

As was intimated in your earlier posting passion is clearly evident. Can this passional orientation lead one astray? You are privy to the full work, as yet unpresented so we will have to wait and see but a question arises in dealing with a pure faith as mentioned with regard to the solitary ascetics: are we to assume that Evola's concept of Initiation is superior to Guenon's? Will we assume that a solitary ascetic(whose practice may reflect either the exterior or the interior-without the discernment of a traditional caste we can never know) can preserve Tradition without any other reference? How are we to make a determination. And is it possible that there has been a continuity so subtle that those of the second function will have difficulty determining its existence in an extant structure without the help of the first function(who may be already embedded within such an expression)?

\hfill

\texttt{Cologero on 2013-02-18 at 15:58 said:}

I don't see where he is saying anything of a sentimental nature or advocating the ``passions'' as somehow superior to intellectuality. On the contrary he writes:

\begin{quote}

Above all, great detachment is necessary, strict intellectuality, the absolute absence of that pseudo-mysticism so much in style at the present time where spirituality and sentimentality are equivalent, where enthusiasm and faith are placed on the same level, where the darkest and murkiest impulsiveness is believed to be an expression of strength, where the external is not only excessive, but tends to destroy the internal, where agenda kills true development, where finally all that is inferior and illegitimate is affirmed with an immodesty which the world has never known before now, even in the most acute periods of decadence.
\end{quote}


The unassailable ``reference point'' would be the ``truths of God'' or Traditional teachings; that is the standard to judge by.

\hfill

\texttt{Tosti on 2013-02-18 at 16:11 said:}

Perhaps it was just the general tone. That may be subjective on my part..my Anglo-saxon temperment.

\hfill

\texttt{Janus on 2013-02-18 at 22:44 said:}

I can see what de Giorgio is saying when he speaks about the central importance of faith. However, I don't understand how he can say that faith is superior to knowledge. He begins by saying (as I understand it) that faith is good because it allows the one without knowledge as yet to move forward in the path, giving him strength and courage. But surely that means that, when knowledge is attained, then faith has fulfilled its purpose? And this seems to imply that knowledge of that which faith led us to pursue is superior, though no less essential, than faith itself.

It also makes me wonder how many priests or ascetics need to arise before a turning of orders can begin. Surely there are some left in Europe, in the old monasteries or as hermits who fulfill this role. But apparently, these few have not prevented the decline, any more than did the Russian hermits in the early 20th century before Bolshevism. Does that mean there was a decline or betrayal among these priests? Who knows... but apparently this function must be fulfilled in a specific function or grow to a certain size before its affects are more generally recognizable. Perhaps this is the distinction between Ascetic and Priests proper? ``Lord, if there is but one righteous man in the city... ?''

\hfill

\texttt{David on 2013-02-18 at 22:57 said:}

@Janus : I cannot answer for De Giorgio, but I don't think we can ever attain perfect knowledge of Reality without being in perfect union with Reality, i.e. God. That is, faith will always be there, because there will always be knowledge. Someone can correct me if I'm wrong.

\hfill

\texttt{Cologero on 2013-02-19 at 00:13 said:}

Janus, please show me where you think De Giorgio says that faith is superior to knowledge... it means I need to make the translation clearer. Here is a quote I have in mind, which says the opposite of the impression you have:

\begin{quotex}

Faith works this miracle that those who cannot reach by an operational effort and aware of the threshold of the divine, attain it with the rapid and direct contact that can be fecund with greater results: compared to those who know, that are beyond the threshold, they find themselves very far, but compared to those who do not believe, to the small men of the small world, the affirmers of the minimum in the minimum and lovers of the shadows, they are in a clearly privileged position because they are distant from just as far as the spirit surpasses the flesh and intelligence surpasses imbecility.
\end{quotex}


Specifically, \emph{compared to those who know, that are beyond the threshold, they find themselves very far} means those who believe are far from those who know, but those who believe are superior to those who neither know nor believe.

As for your other question, let's wait until I get the final part up which addresses some of those issues.

David, perfect knowledge is perfect union as you point out. But I don't understand the other point: where there is knowledge, there is no need for faith.

\hfill

\texttt{Janus on 2013-02-19 at 03:25 said:}

This is the passage:

``Faith is therefore superior to any human knowledge whatsoever, to any conquering activity of the human alone, because it passes beyond it, considers it as ancillary, negligible, negotiable, even nothing in the face of the divine that it recognizes as the invisible venerable root through its invisibility, real through its apparent unreality, divine exactly because not human, not tangible, not discernible with the senses, and analyzable by reason, placed in a sphere of this the good ones, those who believe, will be able to enjoy, if it will be pleasing to God, only after the dissolution of the human compound, i.e., post mortem.''

Perhaps I'm not understanding it correctly. Is he saying here that faith is superior to human knowledge in terms of ``merely human'' ideas about the Divine (it might be useful here to think of the Orthodox doctrine that all human attempts to describe God, while they may be good enough for Man, are ultimately lacking)? I doubt he means the higher Knowledge of gnosis, as that contradicts the passage you mentioned and your point about union and perfect knowledge. The passage you mentioned says what I would have expected, that gnosis is greater than faith, which itself is greater than ignorance.

\hfill

\texttt{Cologero on 2013-02-19 at 08:20 said:}

Not specifically human knowledge about the divine (if there can even be such a thing), but in general. Human knowledge is never more than opinion or a likely story. It does not survive death, as De Giorgio points out, whereas faith certainly does.

\hfill

\texttt{David on 2013-02-19 at 18:47 said:}

My question can be rephrased : how can we attain perfect knowledge if not for perfect being, that is, God, while still here, in the temporal and therefore imperfect realm (on a dualistic gnostic view, or a monistic one) ? I'm just trying to understand, thank you. To me, union with God can be but partial. Or maybe I'm mistaken ?

\hfill

\texttt{Janus on 2013-02-19 at 19:49 said:}

Well, it might help to remember that this unity is in many ways a process of realization of a reality which already is. ``Thou art That''... the singularity of Atman and Brahman. The true Self is the Atman realizing its manifestation and source of being. Thus, a complete awakening, difficult as this might be, should constitute a unity that is total and conscious, and that seems to be whole instead of partial.

\hfill

\texttt{Cologero on 2013-02-20 at 08:11 said:}

David, the rephrasing of the question is better since a response can be split. First of all, there is the question of the possibility of ``perfect knowledge'' or union with God. To that the answer is ``yes'', since what the Intellect knows is ideas or possibilities. That possibility has been described over and over, as there is no shortage of literature. Here on Gornahoor, we try to summarize that literature and re-express it contemporary terms.

Nevertheless, that does not conclude the story, since what you are interested in is the ``actuality'' of the union (theosis), i.e., its realization, not merely its possibility. But that is a question for the Will, not solely the Intellect. Apart from the various claims, there is no scientific way to determine someone's ``state of realization''. We have the testimony of St John Climacos who says in his Ladder of Divine Ascent that theosis is not complete in this life. So perhaps your opinion is aligned with his. Even De Giorgio admits that the steps on the ladder are infinite, so a finite life cannot achieve it.

However, if you are called on that path, that is no excuse to not begin. St Mark the Ascetic in the Philokalia mentions three great enemies:

\begin{enumerate}


\item Ignorance

\item Forgetfulness

\item Sloth
\end{enumerate}


Ignorance is the fist thing to overcome, and obviously Gornahoor readers probably have subdues, or are subduing, that enemy, if they refer to the text we write about. We have addressed the topic of forgetfulness over and over. Never forget who you are, why you were born, what you need to accomplish, what your real life is. For the several hundred daily readers here, we hope, if nothing else, to be their reminders to stay on track, although there are many other resources. Of course, most humans prefer to drink from Lethe, the river in Hell that causes them to forget.

The third enemy is sloth, and in that regard everyone is on his own. No matter how much you know, you cannot expect George\footnote{\url{http://en.wikipedia.org/wiki/Let_George_Do_It_(radio)}} to do it for you. So you will never really know what is possible unless and until the right efforts are made.

\hfill

\texttt{David on 2013-02-20 at 16:58 said:}

That answers it. I will look into St John Climacos to see if I can clarify my mind on the matter. Thank you.

\hfill

\texttt{Saladin on 2013-02-20 at 20:18 said:}

Excellent piece! Loved it! Thanks Cologero! I think De Giorgio is very clear so I don't know why some people are confused. He never says that ``Passions are superior to intellectuality'' (Per Tosti) and neither does he declare the superiority of Faith over TRUE Knowledge (as claimed by Janus). He says Faith is infinitely (Guenon would object to my choice of word here) superior than any worldly (read Philosophical and Scientific) knowledge but at the same time miles away from True Knowledge (i.e. Gnosis). That's what he says, and it's all clear.

I really loved the passage where he lays the blame for degeneration on the Priestly caste. That's a simple fact and needs to be acknowledged. Just as blaming women for Feminism is completely inaccurate: For it is the duty of the Best to insure everything is according to Order (and where the teaching is different, there the Best is lacking!

\hfill

\texttt{Olavsson on 2016-04-05 at 11:24 said:}

Yes, the degeneration of the First Caste of contemplatives is at the origin of the regression of the castes. So if one's true nature, function, purpose and providential destiny (completely without regard for a relative and conditioned, illusory self-will based on forgetfulness) can only be perfectly realized by giving oneself heart and soul, with all one's might, to this path, and even still one chooses to ignore the call in order to invest one's energies in false interests and desires, one is failing in one's true task and responsibility and consequently a serious part of the problem. Too many who potentially (hence not actually) belong to the spiritual caste waste away their lives within profane domains at best producing bourgeois triviality, or at the very worst leading to catastrophic and demonic scientific ``advances.'' My advice can only be: Leave that to the fallen ones---they are many and have no need for you. IF you do in truth have the nature of a contemplative, your true path should already be mapped out clearly: The example to follow is the sages of old, not the profane functions honoured, respected and admired most highly by the modern system. How many more of us can we afford losing to such games falling short of the unworldly super-caste's true principles and agenda? The spiritual conquerors of the contemplative life who uphold the bridge between time and Eternity are fewer with each passing year, whereas the profane classes from high to low have a never-ending supply of new recruits ready to feed the monstrous machine of modernity in whatever way. So it should be obvious enough which task is the most important one for our kind, not only in view of our own spiritual destiny, but in view of the needs of humanity as well.

Indeed, dear De Giorgio is utterly correct regarding Faith. Even a quintessential jnanin-ascetic like Buddha Shakyamuni, founder of a tradition which is as known for its emphasis on Gnosis as any other, many times stressed the great importance of faith, and the superiority of the man possessing unswerving faith as compared to the one who doesn't. Likewise, he listed DOUBT among the most serious obstacles to progress on the Path. In our absurdly degenerate postmodern world, doubt is not seldom glorified almost as if it were a sign of intelligence. But if one lacks the ability to have complete, unfailing and permanent faith as strong as diamond in the truths that are derived from the very Divine Wisdom of the sagely masters, prophets, and avataras, we will never have sufficient determination to go all the way and attain direct and illuminating Knowledge. As far as orthodox essential principles are concerned, starting with sacred metaphysics, there should be no room for doubt if we are serious about our sacred endeavour---such doubt is nothing but a profane weakness hindering absolute resolve on the upward path, and you hardly need complete illumination in order to attain absolute certainty as to the impossibility of every modernist, profane and materialist worldview, upon which the profanely skeptical scientist unbelief-ideology is based, when recognizing its fundamentally flawed principal assumptions. Doubt must be conquered and destroyed long before we can attain complete integral Knowledge. Without perfect faith, we cannot even begin the journey. And without perfect faith in God, in the Supreme, there is only unreality, not even a distant reflection of light.

\texttt{reader on 2013-02-20 at 13:09 said:}

``Those who are called ``geniuses'' by men of today so tied to the human to see them through the greatest irony, also the superhuman... ''

``also the superhuman'' = ``as the superhuman''?

\hfill

\texttt{Cologero on 2013-02-20 at 13:26 said:}

I opted for ``even the superhuman''... the phrasing is a little awkward. If anyone else wants to give it a shot:

Coloro che sono chiamati ``geni'' dagli uomini attuali cosi legati all'umano da scorgervi per suprema ironia, anche il superumano, chiusi come sono nell'ambito umano e terrestre, si trovano di gran lunga inferiori al piu umile dei credenti perche essi ipertrofizzano un nulla, l'uomo e il mondo...

\end{sffamily}
\end{footnotesize}

